\documentclass[UTF8, 12pt, fontset=fandol, oneside]{ctexart}
\usepackage{researchnotes}

\title{凸优化}
\author{}
\date{}

\begin{document}

\maketitle

\section*{1\quad 引言}

引言将对数学优化做一个简单的回顾，我们将主要关注凸优化在数学优化中的特殊地位。对于在引言中出现的一些概念，后面将会更加正式以及详细地进行定义。

\subsection*{1.1\quad 数学优化}

数学优化问题或者说优化问题可以写成如下形式
\begin{equation}
\begin{aligned}
  \text{minimize}\quad & f_0(x)\\
  \text{subject to}\quad & f_i(x)\le b_i,\quad i=1,\cdots,m.
\end{aligned}
\tag{1.1}
\end{equation}
这里，向量 $x=(x_1,\cdots,x_n)$ 称为问题的\textbf{优化变量}，函数 $f_0:\mathbb{R}^n\to\mathbb{R}$ 称为目标函数，函数 $f_i:\mathbb{R}^n\to\mathbb{R}$，$i=1,\cdots,m$，被称为（不等式）\textbf{约束函数}，常数 $b_1,\cdots,b_m$ 称为约束上限或者约束边界。如果在所有满足约束的向量中向量 $x^*$ 对应的目标函数值最小：即对于任意满足约束 $f_1(z)\le b_1,\cdots,f_m(z)\le b_m$ 的向量 $z$，有 $f_0(z)\ge f_0(x^*)$，那么称 $x^*$ 为问题 (1.1) 的\textbf{最优解}或者\textbf{解}。

我们一般考虑具有特殊形式的目标函数和约束函数的优化问题。线性规划即是其中重要的一类。若优化问题 (1.1) 中的目标函数和约束函数 $f_0,\cdots,f_m$ 都是线性函数，即对任意的 $x,y\in\mathbb{R}^n$ 和 $\alpha,\beta\in\mathbb{R}$ 有
\begin{equation}
  f_i(\alpha x+\beta y)=\alpha f_i(x)+\beta f_i(y),
\tag{1.2}
\end{equation}
则此优化问题称为\textbf{线性规划}。若优化问题不是线性的，则称之为\textbf{非线性规划}。

本书讨论一类优化问题，即\textbf{凸优化问题}。凸优化问题中的目标函数和约束函数都是凸函数，即对于任意 $x,y\in\mathbb{R}^n$，任意 $\alpha,\beta\in\mathbb{R}$ 且满足 $\alpha+\beta=1$，$\alpha\ge0$，$\beta\ge0$，下述不等式成立
\begin{equation}
  f_i(\alpha x+\beta y)\le \alpha f_i(x)+\beta f_i(y).
\tag{1.3}
\end{equation}
比较式 (1.3) 和式 (1.2)，可以看出凸性是较线性更为一般的性质：线性函数需要严格满足等式，而凸函数仅仅需要在 $\alpha$ 和 $\beta$ 取特定值的情况下满足不等式。因此线性规划问题也是凸优化问题，可以将凸优化看成是线性规划的扩展。

\subsubsection*{1.1.1\quad 应用}

优化问题 (1.1) 可以看成在向量空间 $\mathbb{R}^n$ 的一集备选解中选择最好的解。用 $x$ 表示备选解，$f_i(x)\le b_i$ 表示 $x$ 必须满足的条件，目标函数 $f_0(x)$ 表示选择 $x$ 的成本（同理也可以认为 $-f_0(x)$ 表示选择 $x$ 的效益或者效用）。优化问题 (1.1) 的解即为满足约束条件的所有备选解中成本最小（或者效用最大）的那个解。

以\textbf{投资组合优化}为例，我们寻求一个最佳投资方案将资本在 $n$ 种资产中进行分配。变量 $x_i$ 表示投资在第 $i$ 种资产上的资本，那么向量 $x\in\mathbb{R}^n$ 则描述了资本在各个资产上的分配情况。约束条件可能包括对投资预算的限制（即总投资额的限制），每一份投资额非负（此处假设不允许空头）以及期望收益必须大于最小可接受收益值。目标函数或者说成本函数可能是对总的风险值或者投资回报方差的一个度量。在此意义下，优化问题 (1.1) 即为在所有可能的投资组合中选择满足所有约束条件且风险最小的那个组合。

另一个例子是电子设计中的器件尺寸问题，即在电子电路中设计每个器件的长度和宽度。此时优化变量表示器件的长度和宽度。约束条件表征了多种工程上需要满足的要求，如生产过程对器件尺寸的要求，电路在特定速度稳定运行对时间的要求，或者是电路的总面积的限制。器件尺寸设计问题中一个比较常见的目标函数是电路的总功耗。因此，优化问题 (1.1) 此时变为设计合适的电路器件尺寸，使之满足设计要求（制造要求、时间要求以及面积要求）并且最为节能。

在\textbf{数据拟合}中，人们需要在一族候选模型中选择最符合观测数据与先验知识的模型。此时，变量为模型中的参数，约束可以是先验知识以及参数限制（比如说非负性）。目标函数可能是与真实模型的偏差或者是观测数据与估计模型的预测值之间的偏差，也有可能是参数值的似然度和置信度的统计估计。优化问题 (1.1) 此时即为寻找合适的模型参数值，使之符合先验知识，且与真实模型之间的偏差或者预测值与观测值之间的偏差最小（或者在统计意义上更加相似）。

大量涉及决策（或系统设计、分析与操作）的实际问题可以表示成数学优化问题，或者数学优化问题的变化形式如多目标优化问题。事实上，数学优化已经成为很多领域的重要工具，它被广泛地应用于工程、电路自动设计、自动控制系统以及土木、化工、机械、航空工程中出现的最优设计问题。此外，网络设计和操作，金融，供应链管理，调度等很多领域的问题都需要用到优化。反映这类应用的列表还在稳定地扩展。

在上述大部分应用中，数学优化都是作为一个辅助工具，协助决策者，系统设计人员或系统操作员根据需要监控系统的运行过程，检查结果以及修改问题（或者方法）。这些人员根据优化问题的结果采取相应的措施，如购买或出售资产以得到最佳投资。

最近出现的一些现象使得数学优化在其他一些领域中的应用成为可能。随着越来越多的计算机嵌入式产品的问世，\textbf{嵌入式优化}的发展势头迅猛。在这些嵌入式应用中，优化的目的是在没有（或者极少）人为干预的条件下实时做出选择甚至实时执行动作。在一些应用领域，传统自动控制系统以及嵌入式优化已经被有机地融合在一起；而在其他一些领域关于这种结合的尝试还处在一个初始阶段。嵌入式实时优化同样带来一些新的挑战：优化结果必须非常可靠且必须在给定的时间（和存储量）内解决问题。

\subsubsection*{1.1.2\quad 求解优化问题}

某类优化问题的\textbf{求解方法}是指（以给定精度）求解此类优化问题中的某一实例的算法。从 20 世纪 40 年代后期开始，学者们致力于设计算法求解多类优化问题，分析这些算法的性质以及进行软件实现。不同算法的有效性，即我们用之求解优化问题 (1.1) 的能力，是大不相同的。它取决于多方面的因素，如目标函数和约束函数的形式，优化问题所包含的变量和约束的个数以及一些特殊的结构，称\textbf{稀疏结构}（如果某个问题中每个约束函数仅仅取决于少数几个变量，那么此问题称为稀疏的）。

即使目标函数和约束函数是光滑的（如多项式），一般形式的优化问题 (1.1) 仍然很难求解。因此，求解一般形式的问题是需要付出一些代价的，如需要较长的计算时间或者可能找不到解。§1.4 讨论了一些这样的算法。

然而，并不是所有的优化问题都难以求解。对于一些特殊的优化问题，存在一些有效的算法，这些算法对含有成百上千变量和约束的大型问题甚至都有效。两类重要且广为人知的例子是最小二乘问题和线性规划问题，随后的 §1.2 将描述这两个问题（第 4 章中将予以详细讨论）。鲜为人知的是还有一类问题，存在有效的求解算法可以如最小二乘以及线性规划问题一样进行有效求解。这类问题即为凸优化问题。一些大型的凸优化问题甚至都能被可靠求解。

\subsection*{1.2\quad 最小二乘和线性规划}

本节介绍广为人知且应用广泛的两类特殊的凸优化问题，最小二乘和线性规划（第 4 章将详细讨论这两类问题）。

\subsubsection*{1.2.1\quad 最小二乘问题}

最小二乘问题是这样一类优化问题，它没有约束条件（即 $m=0$），目标函数是若干项的平方和，每一项具有形式 $a_i^Tx-b_i$，具体形式如下
\begin{equation}
  \text{minimize}\quad f_0(x)=\|Ax-b\|_2^2=\sum_{i=1}^k (a_i^Tx-b_i)^2.
\tag{1.4}
\end{equation}
其中，$A\in\mathbb{R}^{k\times n}$（$k\ge n$），$a_i^T$ 是矩阵 $A$ 的行向量，向量 $x\in\mathbb{R}^n$ 是优化变量。

\noindent\textbf{求解最小二乘问题}

最小二乘问题 (1.4) 的求解可以简化为求解一组线性方程
\[
  (A^TA)x=A^Tb,
\]
因此可得解析解 $x=(A^TA)^{-1}A^Tb$。对于最小二乘问题，有不少好的求解算法（和软件实现），这些算法的精度和可靠性都很高。最小二乘问题可以在有限时间内进行求解，此时间和 $n^2k$ 近似成正比，且比例系数已知。现有的台式计算机可以在几秒之内解决包含数百个变量、数千个求和项的最小二乘问题；当然，更为先进的计算机可以解决更大规模的问题或者以更快的速度解决相同规模的问题（此外，根据摩尔定律，以后的求解时间还会指数下降）。求解最小二乘问题的算法和软件可靠，足以满足嵌入式优化的要求。

很多时候，若最小二乘问题中的系数矩阵 $A$ 具有某些特殊结构，我们甚至可以解决更大规模的此类问题。例如，假设矩阵 $A$ 是稀疏的，即矩阵 $A$ 中含有的非零元素的个数远远少于 $kn$，我们可以更快求解此最小二乘问题，所需时间的数量级远远小于 $n^2k$。现有的台式计算机在一分钟左右的时间内可以处理大规模的稀疏最小二乘问题，这些问题中所含的变量个数甚至多达数万，所含的求和项也可达到数十万（具体时间取决于不同的稀疏类型）。

当最小二乘问题的规模极大（比如说包含上百万的变量）或者需要在给定时间内进行实时求解，此时求解最小二乘问题就有一定的难度。然而，在大部分的应用中，已有的方法还是相当有效和可靠的。事实上，我们可以认为，最小二乘问题的求解是一项（成熟的）\textbf{技术}（除非问题超出目前可解的范围），对于很多人，即使不知道，也无需知道这门技术的细节，仍然可以可靠地应用。

\noindent\textbf{使用最小二乘}

最小二乘问题是回归分析，最优控制以及很多参数估计和数据拟合方法的基础。最小二乘问题有很多统计意义，例如，给定包含高斯噪声的线性测量值时，向量 $x$ 的最大似然估计即等价于最小二乘问题的解。

判断一个优化问题是否是最小二乘问题非常简单：只需要检验目标函数是否是二次函数（然后检验此二次函数是否半正定）。基本最小二乘问题只有一个简单固定的表达式，因此，人们提出了一些标准的技术，使得最小二乘问题在实际应用中更为灵活。

在加权最小二乘问题中，我们最小化加权的最小二乘成本
\[
  \sum_{i=1}^k w_i(a_i^Tx-b_i)^2,
\]
其中，加权系数 $w_1,\cdots,w_k$ 均大于零（加权最小二乘问题可以很方便地转化为标准的最小二乘问题进行求解）。在加权最小二乘问题中，加权系数 $w_i$ 反映了求和项 $a_i^Tx-b_i$ 的重要程度或者说是对解的影响程度。在统计应用中，当给定的线性观测值包含不同方差的噪声时，我们用加权最小二乘来估计向量 $x$。

正则化是解决最小二乘问题的另一个技术，它通过在成本函数中增加一些多余的项来实现。一个最简单的形式是在成本函数中增加一项和变量平方和成正比的项
\[
  \sum_{i=1}^k (a_i^Tx-b_i)^2+\rho\sum_{i=1}^n x_i^2,
\]
这里，$\rho>0$。（此问题亦可以转化为标准最小二乘问题。）当 $x$ 的值较大时，增加的项对其施加一个惩罚，其得到的解比仅优化第一项时更加切合实际。参数 $\rho$ 的选择取决于使用者，选择原则是使原始目标函数值尽可能小而同时保证 $\sum_{i=1}^n x_i^2$ 的值不能太大，在二者之间取得一个较好的平衡。在统计估计中，当待估计向量 $x$ 的分布预先知道时，可以采用正则化方法。

加权最小二乘和正则化在第 6 章会进行详细介绍；第 7 章给出了它们在统计学上的意义。

\subsubsection*{1.2.2\quad 线性规划}

另一类较为重要的优化问题是线性规划，其目标函数和所有的约束函数均为线性函数，线性规划问题可以表述如下
\begin{equation}
\begin{aligned}
  \text{minimize}\quad & c^Tx\\
  \text{subject to}\quad & a_i^Tx\le b_i,\quad i=1,\cdots,m.
\end{aligned}
\tag{1.5}
\end{equation}
其中，向量 $c,a_1,\cdots,a_m\in\mathbb{R}^n$，$b_1,\cdots,b_m\in\mathbb{R}$ 是问题参数，它们决定目标函数和约束函数。

\noindent\textbf{求解线性规划}

虽然线性规划问题的解并没有一个简单的解析表达形式（和最小二乘问题不同），然而，存在很多非常有效的求解线性规划问题的方法，这当中包括 Dantzig 的单纯形法以及最近发展起来的内点法。本书后面将专门介绍内点法。此外，我们也不能给出解决一个线性规划问题所需要的算术运算的确切次数，但对于给定的求解精度，可以给出采用内点法时所需要的算术运算次数的严格上界。其求解复杂度实际正比于 $n^2m$（假设 $m\ge n$），但是比例系数不好确定，而最小二乘方法中比例系数较易确定。尽管与最小二乘问题的算法相比可靠性略逊一等，但求解线性规划问题的算法还是相当可靠的，可以在大约几秒钟的时间内利用小型台式计算机处理含有数百变量、数千约束条件的线性规划问题；如果线性规划问题是稀疏的，或者具有其他有利于运算的结构，甚至可以解决包含数万或者数十万变量的问题。

和最小二乘问题一样，处理极大规模的线性规划问题或者在很短时间内实时解决线性规划问题还是具有一定难度的。但是，和最小二乘问题的情况类似，我们可以说求解（大部分）线性规划问题是一项成熟的技术。线性规划的求解程序可以（并已经）嵌入到很多工具箱和应用软件中。

\noindent\textbf{使用线性规划}

一些应用可以直接表达为线性规划的形式式 (1.5) 或者其他一些标准形式。在很多其他情况中，原始的优化问题并不是线性规划问题的标准形式，但是可以利用第 4 章介绍的技巧转化为一个等价的线性规划问题（然后进行求解）。

作为一个简单例子，考虑 Chebyshev 逼近问题
\begin{equation}
  \text{minimize}\quad \max_{i=1,\cdots,k}|a_i^Tx-b_i|.
\tag{1.6}
\end{equation}
其中，优化变量 $x\in\mathbb{R}^n$，$a_1,\cdots,a_k\in\mathbb{R}^n$。$b_1,\cdots,b_k\in\mathbb{R}$ 为问题的参数，参数取值确定后即得到了一个具体问题。我们注意到，此问题与最小二乘问题式 (1.4) 有着一定的相似性。对于这两个问题，目标函数中均含有 $a_i^Tx-b_i$。不同的是，在最小二乘问题中，我们采用 $a_i^Tx-b_i$ 的平方和作为目标函数，而在 Chebyshev 逼近问题中，我们优化 $a_i^Tx-b_i$ 的绝对值中最大的一项。另外一个重要差别在于 Chebyshev 逼近问题式 (1.6) 中的目标函数是不可微的；而最小二乘问题式 (1.4) 中的目标函数是二次的，自然也是可微的。

求解 Chebyshev 逼近问题 (1.6) 等价于求解如下线性规划问题
\begin{equation}
\begin{aligned}
  \text{minimize}\quad & t\\
  \text{subject to}\quad & a_i^Tx-t\le b_i,\quad i=1,\cdots,k,\\
  & -a_i^Tx-t\le -b_i,\quad i=1,\cdots,k;
\end{aligned}
\tag{1.7}
\end{equation}
优化变量 $x\in\mathbb{R}^n$，$t\in\mathbb{R}$。（具体的细节在第 6 章中进行介绍。）由于线性规划问题的求解非常方便，因此 Chebyshev 逼近问题较易求解。

其实，对于线性规划问题比较熟悉的读者能够很快认识到 Chebyshev 逼近问题式 (1.6) 可以转化为线性规划问题。然而，对于不了解线性规划问题的读者，还是难以将具有不可微目标函数的 Chebyshev 逼近问题式 (1.6) 和线性规划问题联系起来。

尽管判别某个问题是否是可以转化为线性规划问题较之最小二乘问题困难一些，我们仍然可以说这只是一项容易掌握的技术，因为判别是否可以转化为线性规划问题仅仅需要一些标准的技巧。我们甚至可以半自动地完成这种判别转换过程，一些判别以及解决优化问题的软件可以自动识别（一些）可以转化为线性规划的问题。

\subsection*{1.3\quad 凸优化}

凸优化问题具有如下形式
\begin{equation}
\begin{aligned}
  \text{minimize}\quad & f_0(x)\\
  \text{subject to}\quad & f_i(x)\le b_i,\quad i=1,\cdots,m,
\end{aligned}
\tag{1.8}
\end{equation}
其中，函数 $f_0,\cdots,f_m:\mathbb{R}^n\to\mathbb{R}$ 为凸函数。即对于任意 $x,y\in\mathbb{R}^n$，$\alpha,\beta\in\mathbb{R}$ 且 $\alpha+\beta=1$，$\alpha\ge0$，$\beta\ge0$，这些函数满足
\[
  f_i(\alpha x+\beta y)\le \alpha f_i(x)+\beta f_i(y).
\]
从前面的讨论我们可以知道，最小二乘问题 (1.4) 和线性规划问题 (1.5) 实质上都是凸优化问题 (1.8) 的特殊形式。

\subsubsection*{1.3.1\quad 求解凸优化问题}

凸优化问题的解并没有一个解析表达式，但是，和线性规划问题类似，存在很多有效的算法求解凸优化问题。在实际应用中，内点法就较为有效，在一些情况下，可以证明，内点法可以在多项式时间内以给定精度求解这些凸优化问题。（第 11 章将详细讨论这个问题。）

后面将会说明，内点法几乎总可以在 10 步到 100 步之间解决凸优化问题 (1.8)。不考虑特殊结构的凸优化问题（如稀疏结构），每一步需要的操作次数和下述变量成正比
\[
  \max\{n^3,n^2m,F\},
\]
其中 $F$ 是计算目标函数和约束函数 $f_0,\cdots,f_m$ 的一阶导数和二阶导数所需要的计算量。

和线性规划问题类似，这些内点法求解凸优化问题也相当可靠。我们可以使用现在的台式计算机轻易地解决包含数百变量、数千约束的凸优化问题，计算时间不超过一百秒。如果问题本身具有一些特殊结构（如稀疏结构），则可以解决含有数千变量以及约束的更大规模的问题。

然而，求解一般的凸优化问题并不如最小二乘以及线性规划一样是一项成熟的技术。目前，关于内点法在一般非线性凸优化问题中的应用还处于研究阶段，何为最佳求解方法，学者们也看法不一。但是，我们预计，在未来几年内解决一般的凸优化问题将成为一项技术，而这种预计是合理的。事实上，对于凸优化问题的几类重要问题，如二阶锥规划和几何规划问题（第 4 章中将进行详细介绍），内点法正在逐渐成为一项成熟的技术。

\subsubsection*{1.3.2\quad 使用凸优化}

至少从概念上讲，使用凸优化和使用最小二乘以及线性规划类似。如果将某个问题表述为凸优化问题，我们就能迅速有效地进行求解，这点和求解最小二乘问题的情形类似。虽然有点夸张，我们认为，如果某个实际问题可以表述为凸优化问题，那么事实上已经解决了这个问题。

当然，一般的凸优化问题和最小二乘问题以及线性规划问题在某些方面还有很大差异。比如说判断某个问题是否为最小二乘问题非常直接，然而，凸优化问题的识别比较困难。此外，较之线性规划问题，转换为凸优化问题的过程中存在更多的技巧。因此，判断某个问题是否属于凸优化问题或识别那些可以转换为凸优化问题的问题是具有挑战性的工作。本书的主要目的就是为读者建立这方面的知识。当我们具备了识别或者表达一个凸优化问题的技巧时，我们将发现，超乎我们的想象，很多问题都可以利用凸优化求解。

使用凸优化的难点或技巧是判别问题是否属于凸优化问题以及表述问题。一旦实际问题被表述为凸优化问题的形式，解决问题就（几乎）只是一项技术了，就如最小二乘问题和线性规划问题的求解一样。

\subsection*{1.4\quad 非线性优化}

非线性优化（或非线性规划）描述这样一类优化问题，其目标函数或者约束函数是非线性函数，且不一定为凸函数。遗憾的是，对于一般的非线性规划问题式 (1.1)，目前还没有有效的求解方法。有时看似简单的问题，变量个数可能不到 10，却非常难以求解，更不用说上百变量的非线性优化问题。因此，现有的用于求解一般非线性规划问题的方法都是在放宽某些指标的条件下，采取不同的途径进行求解。

\subsubsection*{1.4.1\quad 局部优化}

在\textbf{局部优化}中，人们放宽对解的最优性的要求，不再搜寻使目标函数值最小的最优可行解。取而代之的是，寻找局部最优解。局部最优解是在其小邻域内的所有可行解里使目标函数值最小的那个解，但是不保证优于不在此邻域内的其他可行解。关于一般非线性规划问题的研究有很大一部分关注局部最优化方法，并取得了不少成果。

由于仅仅要求目标函数和约束函数可微，局部优化求解迅速，并可以处理大规模问题。这些特点使得局部优化在实际中的应用较为广泛，尤其是在满意解（非最优解）同样具有价值的应用场合。工程设计就是这种应用的一个例子，此时，局部优化相比经验方法或者其他设计方法已经具有优势，能够改善设计系统的性能。

然而，除了（可能）不能找到全局最优解以外，局部优化方法还存在一些别的缺点。在局部优化中，需要确定优化变量的初始值，而初始值的选取非常重要，对最终得到的局部最优解有着很大的影响。而且，人们无法估计局部最优解相比（全局）最优解到底有多大的差距。此外，局部优化方法对算法的参数值一般也较为敏感，通常需要针对某个具体的问题或某类问题进行调整。

局部优化问题相比最小二乘问题、线性规划问题以及凸优化问题需要更多的技巧。人们需要选择合适的算法，调整算法的参数，选取一个足够好的初始点（针对某个具体问题）或提供一个选取较好初始点的方法（针对一类问题）。大致说来，局部优化方法是一种技巧而不仅是一项技术。局部优化是一种研究得较为透彻的技巧，常常很有效，但是还是一种技巧。相比之下，在最小二乘问题和线性规划问题中需要的技巧很少（当然了，问题规模没有超过目前可以求解的范围）。

比较非线性规划中的局部优化方法和凸优化是一件有意思的事情。大部分局部优化方法仅仅要求目标函数和约束函数可微，因此，将实际问题建模为非线性优化问题是相当直接的。当建模完成后，局部优化中的技巧体现在问题的求解上（在寻找一个局部最优点的意义上）。而凸优化的情形完全相反，技巧和难点体现在描述问题的环节，一旦问题被建模为凸优化问题，求解过程相对来说就非常简单。

\subsubsection*{1.4.2\quad 全局优化}

在\textbf{全局优化}中，人们致力于搜索优化问题 (1.1) 的全局最优解。当然了，付出的代价是效率。在最坏的情况下，全局优化的求解复杂性随着问题的规模 $n$ 和 $m$ 呈指数增长；人们只能寄希望于现实中遇到特殊情况，这样求解速度可能会快很多。虽然这种情况在现实中确实存在，但是并不典型。一般而言，即使问题规模较小，只有几十个变量，求解也需要很长的时间（几个小时甚至几天）。

对变量个数较少的小规模问题，若对计算时间没有苛刻的要求且寻找全局最优解非常有价值，我们采用全局优化。工程设计中高价值系统或安全性第一的系统的最坏情况分析问题或验证问题就是采用全局优化的一个例子。此时，不确定参数是问题的变量，在实际生产过程中或者当工作环境和工作点改变时会发生变化。目标函数是效用函数，函数值越小，情况越坏。关于参数取值的先验知识构成了约束条件。优化问题式 (1.1) 此时即为寻找最坏情况下的参数值，即参数的最差值。如果在最差值的情况下，系统仍然可靠运行，我们认为系统是安全的或者可靠的（当参数发生变化时）。

局部优化方法可以迅速找到一些较差情况下的参数的取值，但是不能保证是最差的情况。如果局部优化方法能够找到一组参数取值，使得系统的性能不在可接受范围内，那么我们说系统是不可靠的。但是，局部优化方法无法证明系统是可靠的，它只是可能没找到最差的情况下的参数取值。与此相反，全局优化能够找到绝对最差的参数取值，如果在此情况下，系统的性能仍然可以接受，我们可以证明系统是安全可靠的。全局优化需要付出的是计算时间，即使对一个参数规模较小的问题都有可能很长。然而，当验证系统可靠非常有价值，或者系统运行不可靠或不安全的代价非常大时，全局优化还是很有必要的。

\subsubsection*{1.4.3\quad 非凸问题中凸优化的应用}

本书主要讨论凸优化问题以及可以转化为凸优化问题的一些应用。不过即使对于非凸问题，凸优化仍然有着重要的作用。

\noindent\textbf{局部优化中利用凸优化进行初始值的选取}

凸优化在非凸问题中的一个重要应用即是将凸优化与局部优化结合起来。对于一个非凸问题，我们首先将其表述为近似凸优化问题。我们通过求解近似凸问题，得到近似问题的精确解。这是很容易做到的，因为凸优化的求解不需要初始值，求解容易。然后，我们用凸问题的精确解作为局部优化算法的初始值，求解原非凸问题。

\noindent\textbf{非凸优化中的凸启发式算法}

求解非凸问题的很多启发式算法都是基于凸优化。搜寻满足一定约束条件的稀疏向量就是这方面应用的一个有趣的例子，前面已经说过，稀疏向量即含有较少非零元素的向量。此问题是一个复杂的组合问题，然而，基于凸优化，存在一些简单的启发式算法，可以找到较稀疏的解。（第 6 章将对此进行详细介绍。）

另一个较为常见的应用例子是随机算法。随机算法产生服从某个概率分布的一些备选解，在这些备选解中选择最好的那个解作为非凸问题的近似解。假设产生备选解的这些概率分布可以由参数表征，比如它的均值和方差。我们就可以提出这样的问题，在所有这些分布中，哪个分布使得目标函数具有最小的期望？事实证明，这个问题有时可以表述为凸问题，因此可以被有效求解（例如可以参见习题 11.23）。

\noindent\textbf{全局优化的界}

对于非凸问题的全局优化，很多方法都需要给出最优解的下界，而且计算代价必须较小。求解下界的两个标准方法都是基于凸优化。在松弛算法中，每个非凸约束都用一个松弛的凸约束代替。在 Lagrange 松弛中，我们求解 Lagrange 对偶问题（在第 5 章介绍）。此问题是凸的，并给出了原非凸问题最优解的一个下界。

\subsection*{1.5\quad 本书主要内容}

本书分为三个主要部分，分别为理论、应用以及算法。

\subsubsection*{1.5.1\quad 第 I 部分：理论}

在第 I 部分理论中，我们介绍基本的定义、概念以及凸分析和凸优化的一些结果。我们不力求面面俱到，只是对我们认为在判别、表达凸优化问题中有用的理论进行详述。这些都属于经典问题，几乎所有问题都可以在其他有关凸分析和凸优化的教材中找到。本书并不打算给出这些问题的最一般性的结果，读者如果对此有兴趣可以参照其他有关凸分析的标准教材。

第 2 章和第 3 章分别介绍凸集和凸函数，这两章将给出几种常见的凸集和凸函数。此外，我们还给出几种凸运算规则，即针对凸集以及凸函数的保凸运算。结合常见的例子以及凸运算规则，我们可以得出（或者更重要的，判别）一些比较复杂的凸集和凸函数。

在第 4 章凸优化问题中，我们对优化问题进行详细的讨论，并且介绍几种变换用来将问题表述为凸优化问题。此外，我们还介绍几类常用的凸优化问题，如线性规划、几何规划以及最近发展起来的二阶锥规划和半定规划。

第 5 章对偶讲述了在凸优化求解中至关重要的 Lagrange 对偶理论，给出了经典的 Karush-Kuhn-Tucker 最优性条件以及凸优化问题的局部敏感性和全局敏感性分析。

\subsubsection*{1.5.2\quad 第 II 部分：应用}

在第 II 部分应用中，我们介绍凸优化在概率统计、计算几何以及数据拟合中的广泛应用。我们力求给出直观的描述，使得大部分读者都能够理解这些应用。为了使应用部分不致太长，我们仅考虑简单的例子，对于一些问题也会增加可能的扩展讨论。这种处理问题的方式可能会让一些专家认为不够专业，对此，我们表示歉意。但是，本书的目的是传达应用的精髓，使得更多读者能够更快接受，而不是从理论上详细、完整地进行介绍。我们自己的专业背景是电气工程，涉及的领域是控制系统、信号处理和电路分析设计。虽然在我们的课堂上（以本书作为主要教材）大量提到我们专业方面的应用，但是我们在书中仅仅讨论了这些方面的一些应用。

第 II 部分的目的是让读者通过例子了解如何在实践中应用凸优化。

\subsubsection*{1.5.3\quad 第 III 部分：算法}

在第 III 部分算法中，我们介绍了求解凸优化问题的数值方法，主要介绍牛顿算法和内点法。第 III 部分包括 3 章，分别涉及无约束优化、等式约束优化以及不等式约束优化问题。这三个章节按照问题由易到难的顺序进行安排，排在后面章节中的问题总可以表述为求解一系列前面章节所涉及的较简单的问题。二次规划问题（包含最小二乘问题）是最简单的问题，可以通过求解线性方程组进行有效求解。第 9 章中所介绍的牛顿法求解次简单的问题。在牛顿法中，需要求解的问题为无约束或者等式约束的问题，可以将这些问题转化为求解一系列二次规划问题。第 11 章介绍求解最难问题的内点法。这些方法通过求解一系列的无约束或者等式约束问题来求解不等式约束问题。

总而言之，我们仅仅根据需要介绍几种算法，对于另外一些好方法，如拟牛顿法、共轭梯度法、bundle 方法以及割平面方法，本书并未涉及。对于所介绍的方法，本书给出简化的形式，关于其最新的、最复杂的形式都未涉及。在选择介绍什么算法时，我们考虑了多方面的因素。本文所介绍的算法都较为简单（无论描述或是实现），对于大多数问题不但快速有效，而且鲁棒可靠。

有不少用户都是仅仅使用（并不开发）凸优化软件，如线性规划或半定规划求解软件。对于这些用户，第 III 部分所覆盖的知识旨在传达凸优化方法的基本概念和属性。对于一些开发新算法的用户，我们相信，第 III 部分提供了一个较好的初步介绍。

\subsubsection*{1.5.4\quad 附录}

本书有三个附录。第一个附录列举了本书可能用到的数学知识并定义了本书所使用的数学符号。第二个附录讨论一类特殊的优化问题，其目标函数是二次函数且含有一个二次约束。这类问题是非凸问题，但可以利用第 II 部分一些应用中的结果进行有效求解。

最后一个附录简要介绍了数值线性代数，主要关注一些方法。这些方法可以利用问题结构，如稀疏性，从而使得问题的求解更为有效。我们并未涉及一些重要的话题，如舍入分析，也没有给出实现所需要的因式分解的方法的技术细节。事实上，这些问题在很多优秀的教材中都有涉及。

\subsubsection*{1.5.5\quad 关于例子}

本书的很多地方（尤其是在介绍应用和算法的第 II 和第 III 部分）都利用具体的例子方便读者更好地理解本书的方法。有的时候，我们选择（或设计）例子来解释我们的观点；其他时候，我们选择“典型”例子。所谓典型例子是指利用一些简单明了的概率分布随机产生的例子。我们知道，仅仅用几十个或者上百个随机产生的例子去评价某个算法的性能是不可靠的，所以本书并非用例子来精确估计算法性能。我们给出这些例子只是希望读者对算法的性能，或者问题规模对计算量的影响有个大致的了解。事实上，对同样的例子，读者的求解结果可能和我们的不一样。

\subsubsection*{1.5.6\quad 关于习题}

每一章的结尾都有一些习题。有些习题针对正文的某些观点展开详细的讨论；有些习题则侧重于判断或建立给定集合、函数或者问题的凸性；另外还有一些习题，考虑如何构造表述一个凸优化问题。有些章节包含数值习题，可能需要用合适的高级语言进行一些编程（但是这种情形不多）。习题的难度不一，有简单明了的也有需要一些技巧的。这些习题混杂在一起，其难度均未标出。

\subsection*{1.6\quad 符号}

除了一些特例，本书所采用的基本是标准符号，本节介绍基本符号；更详细的符号列表见书后符号表。

实数集用 $\mathbb{R}$ 表示，$\mathbb{R}_+$ 表示非负实数的集合，而 $\mathbb{R}_{++}$ 表示正实数的集合。$\mathbb{R}^n$ 表示 $n$ 维实向量的集合，$\mathbb{R}^{m\times n}$ 表示 $m\times n$ 实矩阵的集合。在本书中，向量和矩阵均用方括号限定，不同的元素用空格区分。若向量元素之间用逗号区分并包含在圆括号内，其代表列向量。例如，如果 $a,b,c\in\mathbb{R}$，对于 $\mathbb{R}^3$ 中的某个向量，我们有
\[
  (a,b,c)=
  \begin{bmatrix}
  a\\
  b\\
  c
  \end{bmatrix}
  =
  \begin{bmatrix}
  a & b & c
  \end{bmatrix}^T.
\]
符号 $\mathbf{1}$ 表示所有分量均为 1 的向量（维数由实际情况决定）。$x_i$ 可以表示向量 $x$ 的第 $i$ 个分量，也可以表示序列或者向量集 $x_1,x_2,\cdots$ 的第 $i$ 个元素。具体用到这些符号的时候，根据上下文可以知道它们的涵义。

$\mathbf{S}^k$ 表示 $k\times k$ 对称矩阵的集合，$\mathbf{S}_+^k$ 表示 $k\times k$ 对称半正定矩阵的集合，而 $\mathbf{S}_{++}^k$ 表示 $k\times k$ 对称正定矩阵的集合。弯不等式符号 $\succeq$（以及其严格形式 $\succ$）用来表示扩展的不等式：对向量而言，它刻画了每个分量之间的不等式关系；对于对称矩阵，它表征了矩阵不等式。若添加下标，符号 $\succeq_K$（或者 $\succ_K$）表示相对于锥 $K$ 的广义不等式（将在 §2.4.1 中详细说明）。

在本书中，描述函数的方式同标准符号有点不同，希望不会引起误解。我们使用符号 $f:\mathbb{R}^p\to\mathbb{R}^q$ 表示函数 $f$ 定义在 $\mathbb{R}^p$ 的某个子集上，且输出值在 $\mathbb{R}^q$ 空间内。这个子集即为函数 $f$ 的定义域，用符号 $\mathbf{dom}\,f$ 描述。可以认为符号 $f:\mathbb{R}^p\to\mathbb{R}^q$ 代表了函数类型，因为在计算机语言中，$f:\mathbb{R}^p\to\mathbb{R}^q$ 意味着函数 $f$ 的输入是一个 $p$ 维的实向量，输出是一个 $q$ 维的实向量。在本书中，$f$ 的定义域 $\mathbf{dom}\,f$ 是 $\mathbb{R}^p$ 中的子集，$x$ 必须是这个子集中的点。比如说对数函数 $\log:\mathbb{R}\to\mathbb{R}$，其定义域为 $\mathbf{dom}\,\log=\mathbb{R}_{++}$。符号 $\log:\mathbb{R}\to\mathbb{R}$ 意味着对数函数的输入和输出均为实数；而 $\mathbf{dom}\,\log=\mathbb{R}_{++}$ 说明对数函数仅仅对大于零的实数有意义。

我们使用 $\mathbb{R}^n$ 表示一般的有限维空间。在本书中，还有其他一些有限维空间，如最高次数给定的单变量多项式空间，或者 $k\times k$ 对称矩阵空间 $\mathbf{S}^k$。通过选择某个向量空间的基，可以将其与 $\mathbb{R}^n$ 联系起来（$n$ 是其维数），此时关于向量空间 $\mathbb{R}^n$ 的有关结论都可以应用。至于如何建立联系并转换结论至其他向量空间，一般留给读者完成。例如，任意线性函数 $f:\mathbb{R}^n\to\mathbb{R}$ 可以表示为 $f(x)=c^Tx$，其中 $c\in\mathbb{R}^n$。相应地关于向量空间 $\mathbf{S}^k$ 的描述可以通过选定基并进行转换而得到，即：任意线性函数 $f:\mathbf{S}^k\to\mathbb{R}$ 都可以表征为 $f(X)=\mathbf{tr}(CX)$，其中 $C\in\mathbf{S}^k$。

\section*{参考文献}

最小二乘是一个古老的课题；Gauss 在 19 世纪 20 年代撰写的论文（拉丁文）就曾经提到过，最近由 Stewart [Gau95] 进行了翻译。更近的工作可以参看 Lawson 和 Hanson [LH95] 以及 Björck [Bjö96] 所撰写的书。关于线性规划的参考文献将在第 4 章中提到。

关于非线性规划问题，有一些很好的求解局部最优解的算法，这些算法可以在下列文献中找到，Gill, Murray 和 Wright [GMW81]，Nocedal 和 Wright [NW99]，Luenberger [Lue84]，以及 Bertsekas [Ber99]。

下列书籍讲述了全局优化，Horst 和 Pardalos [HP94]，Pinter [Pin95]，以及 Tuy [Tuy98]。利用凸优化问题寻找非凸问题的界是一个较热的研究课题；在上述全局优化的书中，Ben-Tal 和 Nemirovski 撰写的书 [BTN01, §4.3] 中以及 Nesterov，Wolkowicz 和 Ye 写的综述 [NWY00] 均有提及。关于此课题的一些重要的文章可以参看 Goemans 和 Williamson [GW95]，Nesterov [Nes00, Nes98]，Ye [Ye99]，以及 Parrilo [Par03]。关于随机方法，Motwani 和 Raghavan 的文章中进行了讨论 [MR95]。

凸分析，即关于凸集，凸函数以及凸优化问题的数学分析，是数学问题的一个较为成熟的分支。基本的参考文献可以参看书籍 Rockafellar [Roc70]，Hiriart-Urruty 和 Lemaréchal [HUL93, HUL01]，Lemaréchal [HUL93, HUL01]，Borwein 和 Lewis [BL00]，以及 Bertsekas，Nedić 和 Ozdaglar [Ber03]。更多的关于凸分析的参考文献列在第 2 章至第 5 章中。

Nesterov 和 Nemirovski [NN94] 首先提出内点法可以解决很多凸优化问题；类似的文献在第 11 章中列出。关于现代凸优化问题，内点法及其应用可以参看 Ben-Tal 和 Nemirovski 所写的书 [BTN01]。

另外还有一些方法用来求解凸优化问题，本书没有涉及，如次梯度法 [Sho85]，bundle 方法 [HUL93]，割平面法 [Kel60, EM75, GLY96] 以及椭球法 [Sho91, BGT81] 等。

凸优化问题比较容易求解的思想由来已久。很早以前，人们就认为凸优化的理论相对一般的非线性优化理论直接（完备）得多。下面是 Rockafellar 在 1993 年 SIAM 的综述中曾经写过的一段话 [Roc93]：

“事实上，优化问题的分水岭不是线性和非线性，而是凸性和非凸性。”

Nemirovski 和 Yudin，在他们的著作优化中的问题复杂性以及方法效率 [NY83] 中第一次正式提出凸优化问题较之一般的非线性优化问题更易求解。他们说明，凸优化问题的信息复杂度要远远低于一般的非线性优化问题。更新的关于此类话题的著作可以参看 Vavasis [Vav91]。内点法具有较低的复杂度（理论上）和现代的相关研究不谋而合。现代很多研究致力于证明内点法（或其他方法）能够在有限步运算内求解一些凸优化问题，而所谓的有限步是指增长速度不超过问题维数和 $\log(1/\epsilon)$ 的多项式，其中 $\epsilon>0$ 是要求的精度（具体可以参看第 11 章中的一些简单结果）。第一个全面研究此问题的著作是 Nesterov 和 Nemirovski [NN94]。其他的参考文献可以参看 Ben-Tal 和 Nemirovski [BTN01, 第 5 讲] 以及 Renegar [Ren01]。在求解很多凸优化问题时内点法具有多项式时间的复杂度，这大大优于利用现有算法求解许多非凸优化问题的情形，事实上，在最坏的情况下，这些算法需要的运算次数随着问题的维数量指数增长。

凸优化的应用非常广泛，在此难以一一罗列。凸分析在经济和金融领域构成了很多结论的基础，具有至关重要的地位。比如说，在无套利的假设条件下，分离超平面定理可以用来推导价格和风险中性概率的存在性（可以参看 Luenberger [Lue95, Lue98] 和 Ross [Ros99]）。凸优化，尤其是我们可以解决半定规划的能力，在自动控制理论上也得到了广泛的关注。凸优化在控制理论中的应用可以参看下列著作，Boyd 和 Barrat [BB91]，Boyd，El Ghaoui，Feron 和 Balakrishnan [BEFB94]，Dahleh 和 Diaz-Bobillo [DDB95]，El Ghaoui 和 Niculescu [EN00]，以及 Dullerud 和 Paganini [DP00]。另一个嵌入式（凸）优化问题是预测控制。预测控制是需要在每一步求解一个二次规划（凸优化）问题的自动控制策略。预测控制如今已在化学工业的过程控制中得到了广泛的应用；见 Morari 和 Zafiriou [MZ89]。凸优化方法（尤其是几何规划方法）的另一个应用领域是电子电路设计，在此领域的研究历史悠久。这方面的文章有 Fishburn 和 Dunlop [FD85]，Sapatnekar，Rao，Vaidya 和 Kang [SRVK93]，以及 Hershenson，Boyd 和 Lee [HBL01]。Luo 在文 [Luo03] 中对凸优化在信号处理以及通信方面的应用做了一个综述。更多的关于凸优化的应用的参考文献可以在第 4 章以及第 6 章和第 8 章中找到。

最近用来求解凸优化问题的内点法的高效实现可以参考软件包 LOQO [Van97] 和 MOSEK [MOS02] 以及第 11 章中列出的代码。详细描述优化问题的软件系统有 AMPL [FGK99] 和 GAMS [BKMR98]。这两个软件系统都提供了一些方法，判别可以转化为线性规划的问题。

\end{document}
